% GENERATED FILE. Edit canonical lessons, then run npm run generate:books.
% canonical-source-hash: fnv1a64:755dcec8664db150
% canonical-lessons: LA-C541-balneator, LA-C541-caldarium, LA-C541-tepidarium, LA-C541-frigidarium, LA-C541-piscina

\chapter{At the baths}
\label{ch:la-at-the-baths}
\hlchaptermodality{541}

\begin{chapteropening}
\textbf{By the end of this chapter:} \emph{I can name the rooms of the Roman baths.}

Complete the last lesson of chapter 541: I can name the rooms of the Roman baths.
\end{chapteropening}

\section[balneātor, balneātōris]{balneātor, balneātōris — a bath-keeper}

\label{lesson:LA-C541-balneator}



\begin{quote}
Before the new one: say the Latin for a hand-ball game, then the Latin for swimming.
\end{quote}

\subsection*{You'll want to know: balneātor, balneātōris}
\textbf{balneātor, balneātōris} — \textquotedblleft{}a bath-keeper\textquotedblright{}.

\begin{quote}
Balneātor aquam calefacit.
\end{quote}

\begin{quote}
\emph{The bath-keeper heats the water.}
\end{quote}

Masculine, third declension, from \emph{balneum}: the man who runs a bath-house.

\begin{cousinweb}
\textbf{balneātor} is built on \emph{balneum} \textquotedblleft{}bath\textquotedblright{} with \emph{-tor} (Lewis \& Short); \emph{balneum} was borrowed from Greek \emph{balaneion}.
\end{cousinweb}

\begin{grammarlens}[title={What you've built}]
One more word for telling, explaining and describing at B1.
\end{grammarlens}

\subsection*{Your turn}
\begin{itemize}
\raggedright
  \item \emph{Say it:} \emph{balneātor, balneātōris}
  \item \emph{Say it:} \emph{balneātor, balneātōris}, once more
  \item \emph{Say it:} a sentence of your own with \emph{balneātor, balneātōris}
  \item \emph{Recall:} say the Latin for a hand-ball game, then the Latin for swimming, then say \emph{balneātor, balneātōris} again
\end{itemize}

\begin{tcolorbox}[breakable,colback=teal!4,colframe=teal!35!black,title={Before you move on}]
What does balneātor, balneātōris mean? (\textquotedblleft{}A bath-keeper\textquotedblright{}.) Say it once more.
\end{tcolorbox}
\section[caldārium, caldāriī]{caldārium, caldāriī — the hot room (of the baths)}

\label{lesson:LA-C541-caldarium}



\begin{quote}
Before the new one: say the Latin for swimming, then the Latin for a bath-keeper.
\end{quote}

\subsection*{You'll want to know: caldārium, caldāriī}
\textbf{caldārium, caldāriī} — \textquotedblleft{}the hot room (of the baths)\textquotedblright{}.

\begin{quote}
In caldāriō aqua calida est.
\end{quote}

\begin{quote}
\emph{In the hot room the water is hot.}
\end{quote}

Neuter, second declension, from \emph{caldus} — a short form of \emph{calidus}, hot.

\begin{cousinweb}
\textbf{caldārium} is built on \emph{caldus}, a shortened \emph{calidus} \textquotedblleft{}hot\textquotedblright{} (Lewis \& Short; EtymDB). English \emph{cauldron} comes from the related \emph{caldāria} \textquotedblleft{}cooking pot\textquotedblright{}.
\end{cousinweb}

\begin{grammarlens}[title={What you've built}]
One more word for telling, explaining and describing at B1.
\end{grammarlens}

\subsection*{Your turn}
\begin{itemize}
\raggedright
  \item \emph{Say it:} \emph{caldārium, caldāriī}
  \item \emph{Say it:} \emph{caldārium, caldāriī}, once more
  \item \emph{Say it:} a sentence of your own with \emph{caldārium, caldāriī}
  \item \emph{Recall:} say the Latin for swimming, then the Latin for a bath-keeper, then say \emph{caldārium, caldāriī} again
\end{itemize}

\begin{tcolorbox}[breakable,colback=teal!4,colframe=teal!35!black,title={Before you move on}]
What does caldārium, caldāriī mean? (\textquotedblleft{}The hot room (of the baths)\textquotedblright{}.) Say it once more.
\end{tcolorbox}
\section[tepidārium, tepidāriī]{tepidārium, tepidāriī — the warm room}

\label{lesson:LA-C541-tepidarium}



\begin{quote}
Before the new one: say the Latin for a bath-keeper, then the Latin for the hot room.
\end{quote}

\subsection*{You'll want to know: tepidārium, tepidāriī}
\textbf{tepidārium, tepidāriī} — \textquotedblleft{}the warm room\textquotedblright{}.

\begin{quote}
Ē caldāriō in tepidārium īmus.
\end{quote}

\begin{quote}
\emph{From the hot room we go into the warm room.}
\end{quote}

Neuter, second declension, from \emph{tepidus} — lukewarm, already known.

\begin{cousinweb}
\textbf{tepidārium} is built on \emph{tepidus} \textquotedblleft{}lukewarm\textquotedblright{} (Lewis \& Short).
\end{cousinweb}

\begin{grammarlens}[title={What you've built}]
One more word for telling, explaining and describing at B1.
\end{grammarlens}

\subsection*{Your turn}
\begin{itemize}
\raggedright
  \item \emph{Say it:} \emph{tepidārium, tepidāriī}
  \item \emph{Say it:} \emph{tepidārium, tepidāriī}, once more
  \item \emph{Say it:} a sentence of your own with \emph{tepidārium, tepidāriī}
  \item \emph{Recall:} say the Latin for a bath-keeper, then the Latin for the hot room, then say \emph{tepidārium, tepidāriī} again
\end{itemize}

\begin{tcolorbox}[breakable,colback=teal!4,colframe=teal!35!black,title={Before you move on}]
What does tepidārium, tepidāriī mean? (\textquotedblleft{}The warm room\textquotedblright{}.) Say it once more.
\end{tcolorbox}
\section[frīgidārium, frīgidāriī]{frīgidārium, frīgidāriī — the cold room}

\label{lesson:LA-C541-frigidarium}



\begin{quote}
Before the new one: say the Latin for the hot room, then the Latin for the warm room.
\end{quote}

\subsection*{You'll want to know: frīgidārium, frīgidāriī}
\textbf{frīgidārium, frīgidāriī} — \textquotedblleft{}the cold room\textquotedblright{}.

\begin{quote}
Post caldārium in frīgidāriō natāmus.
\end{quote}

\begin{quote}
\emph{After the hot room we swim in the cold room.}
\end{quote}

Neuter, second declension, from \emph{frīgidus} — cold. It held a cold plunge pool.

\begin{cousinweb}
\textbf{frīgidārium} is built on \emph{frīgidus} \textquotedblleft{}cold\textquotedblright{} (Lewis \& Short); EtymDB traces \emph{frīgus} to a Proto-Indo-European word for \textquotedblleft{}cold\textquotedblright{}. English \emph{frigid} comes from \emph{frīgidus}.
\end{cousinweb}

\begin{grammarlens}[title={What you've built}]
One more word for telling, explaining and describing at B1.
\end{grammarlens}

\subsection*{Your turn}
\begin{itemize}
\raggedright
  \item \emph{Say it:} \emph{frīgidārium, frīgidāriī}
  \item \emph{Say it:} \emph{frīgidārium, frīgidāriī}, once more
  \item \emph{Say it:} a sentence of your own with \emph{frīgidārium, frīgidāriī}
  \item \emph{Recall:} say the Latin for the hot room, then the Latin for the warm room, then say \emph{frīgidārium, frīgidāriī} again
\end{itemize}

\begin{tcolorbox}[breakable,colback=teal!4,colframe=teal!35!black,title={Before you move on}]
What does frīgidārium, frīgidāriī mean? (\textquotedblleft{}The cold room\textquotedblright{}.) Say it once more.
\end{tcolorbox}
\section[piscīna, piscīnae]{piscīna, piscīnae — a pool; a fish-pond}

\label{lesson:LA-C541-piscina}



\begin{quote}
Before the new one: say the Latin for the warm room, then the Latin for the cold room.
\end{quote}

\subsection*{You'll want to know: piscīna, piscīnae}
\textbf{piscīna, piscīnae} — \textquotedblleft{}a pool; a fish-pond\textquotedblright{}.

\begin{quote}
In piscīnā thermārum puerī natant.
\end{quote}

\begin{quote}
\emph{Boys swim in the pool of the baths.}
\end{quote}

Feminine, first declension, from \emph{piscis} — fish: first a fish-pond, then any pool for swimming.

\begin{cousinweb}
\textbf{piscīna} is built on \emph{piscis} \textquotedblleft{}fish\textquotedblright{} (Lewis \& Short; EtymDB): a fish-pond, then a swimming pool. French \emph{piscine} \textquotedblleft{}swimming pool\textquotedblright{} comes from it.
\end{cousinweb}

\begin{grammarlens}[title={What you've built}]
One more word for telling, explaining and describing at B1.
\end{grammarlens}

\subsection*{Your turn}
\begin{itemize}
\raggedright
  \item \emph{Say it:} \emph{piscīna, piscīnae}
  \item \emph{Say it:} \emph{piscīna, piscīnae}, once more
  \item \emph{Say it:} a sentence of your own with \emph{piscīna, piscīnae}
  \item \emph{Recall:} say the Latin for the warm room, then the Latin for the cold room, then say \emph{piscīna, piscīnae} again
\end{itemize}

\begin{tcolorbox}[breakable,colback=teal!4,colframe=teal!35!black,title={Before you move on}]
What does piscīna, piscīnae mean? (\textquotedblleft{}A pool; a fish-pond\textquotedblright{}.) Say it once more.
\end{tcolorbox}
